\section{Paired Gamma deformations and polygamma integrals}
\label{gn:sec:paired}

The parameter freedom in Theorem~\ref{gn:thm:determinant} is useful
beyond the single quarter specialization.  For $n\ge1$, $\sigma>0$,
and $m=n+1$, take
\begin{equation}\label{gn:eq:paired-parameters}
 a_{2r-1}=\sigma+u_r,\qquad a_{2r}=\sigma-u_r,
 \qquad 1\le r\le m,\qquad |\Re u_r|<\sigma.
\end{equation}
Let $\mathcal D_{n,\sigma}(\mathbf u)$ denote the corresponding
determinant.  Every pair has the same sum $2\sigma$, so the denominator
in \eqref{gn:eq:determinant} remains fixed.

\begin{theorem}[Complete paired deformation]
\label{gn:thm:paired}
On the domain \eqref{gn:eq:paired-parameters},
\begin{equation}\label{gn:eq:paired-gamma}
 \boxed{\displaystyle
 \frac{\mathcal D_{n,\sigma}(\mathbf u)}
      {\mathcal D_{n,\sigma}(\mathbf0)}
 =\prod_{r<s}
   \frac{\displaystyle\prod_{\varepsilon,\eta\in\{-1,1\}}
          \Gamma(2\sigma+\varepsilon u_r+\eta u_s)}
        {\Gamma(2\sigma)^4}.}
\end{equation}
At the quarter parameter this is elementary:
\begin{equation}\label{gn:eq:paired-sec}
 \boxed{\displaystyle
 \frac{\mathcal D_{n,1/4}(\mathbf u)}
      {\mathcal D_{n,1/4}(\mathbf0)}
 =\prod_{r<s}\sec\bigl(\pi(u_r+u_s)\bigr)
                 \sec\bigl(\pi(u_r-u_s)\bigr).}
\end{equation}
If only one pair is deformed, the two respective ratios are
\begin{equation}\label{gn:eq:one-pair}
 \left\{\frac{\Gamma(2\sigma+u)\Gamma(2\sigma-u)}
                    {\Gamma(2\sigma)^2}\right\}^{2n},
 \qquad \sec^{2n}(\pi u).
\end{equation}
\end{theorem}

\begin{proof}
The $m$ within-pair Gamma factors are $\Gamma(2\sigma)$ and are
unchanged.  Each pair of distinct labels contributes the four factors
shown in \eqref{gn:eq:paired-gamma}; all remaining constants cancel.
At $\sigma=1/4$, apply
$\Gamma(1/2+z)\Gamma(1/2-z)=\pi/\cos(\pi z)$ first to
$z=u_r+u_s$ and then to $z=u_r-u_s$.  The single-pair specialization
has $n$ identical pair interactions and gives \eqref{gn:eq:one-pair}.
\end{proof}

\subsection{A positive density and its exact normalization}

Define
\begin{equation}\label{gn:eq:density}
\begin{split}
 \rho_{n,\sigma}(\mathbf x)
 ={}&\Delta(x_1^2,\ldots,x_n^2)
      \Delta(\cosh\pi x_1,\ldots,\cosh\pi x_n)\\
 &\times\prod_{j=1}^n x_j\sinh(\pi x_j)
              |\Gamma(\sigma+ix_j/2)|^{4(n+1)},
 \qquad \mathbf x\in(0,\infty)^n.
\end{split}
\end{equation}
The density is symmetric and nonnegative, and is positive off the
collision hyperplanes.  Its normalizing integral is
\begin{equation}\label{gn:eq:Z}
 \mathcal Z_{n,\sigma}
 =2^{n(3n+1)/2}\pi^{n(n+1)}n!
       \frac{\Gamma(2\sigma)^{(n+1)(2n+1)}}
            {\Gamma(2(n+1)\sigma)}.
\end{equation}
This follows from \eqref{gn:eq:I-evaluation}.  Write $\mathbb E_{n,\sigma}$
for expectation with respect to $\rho_{n,\sigma}/\mathcal Z_{n,\sigma}$,
and define the real linear statistics
\begin{equation}\label{gn:eq:statistics}
 S_q=\sum_{j=1}^n\Re\psi^{(q)}(\sigma+ix_j/2),
                   \qquad q=1,3,5,\ldots.
\end{equation}

\begin{theorem}[All-order polygamma insertion identity]
\label{gn:thm:polygamma-germ}
As an identity of holomorphic germs at $\mathbf u=\mathbf0$,
\begin{equation}\label{gn:eq:polygamma-germ}
\begin{split}
 &\mathbb E_{n,\sigma}
  \exp\left(4\sum_{r=1}^{n+1}\sum_{k\ge1}
                   \frac{u_r^{2k}}{(2k)!}S_{2k-1}\right)\\
 &\hspace{12mm}=
 \prod_{r<s}
   \frac{\displaystyle\prod_{\varepsilon,\eta\in\{-1,1\}}
         \Gamma(2\sigma+\varepsilon u_r+\eta u_s)}
        {\Gamma(2\sigma)^4}.
\end{split}
\end{equation}
In particular, every finite coefficient gives an ordinary absolutely
convergent integral identity in polygamma functions.
\end{theorem}

\begin{proof}
Apply continuous Cauchy--Binet to the deformed determinant and divide
by its value at zero.  The interaction Vandermondes are unchanged, so
the quotient is the expectation of the product of Gamma-factor ratios.
For one pair and $z=\sigma+ix/2$, the logarithm of that ratio is
\[
\begin{aligned}
 &\log\Gamma(z+u)+\log\Gamma(z-u)-2\log\Gamma(z)\\
 &\quad+\log\Gamma(\overline z+u)+\log\Gamma(\overline z-u)
                         -2\log\Gamma(\overline z)\\
 &\hspace{18mm}=4\sum_{k\ge1}\frac{u^{2k}}{(2k)!}
                                      \Re\psi^{(2k-1)}(z).
\end{aligned}
\]
The logarithmic Gamma branches are holomorphic on the right half-plane.
Their Taylor expansions converge uniformly for $|u|$ in a fixed
sufficiently small disk, with the integral dominated as in
Theorem~\ref{gn:thm:determinant}.  Alternatively every fixed coefficient
can be obtained directly by differentiating the finite Gamma product
under its normally convergent integral.  Multiplication over spatial
variables and pair labels gives the left side of
\eqref{gn:eq:polygamma-germ}; Theorem~\ref{gn:thm:paired} evaluates it.
\end{proof}

For instance, the single-pair coefficients at degrees two and four are
\begin{align}
 \mathbb E_{n,\sigma}S_1
      &=n\psi'(2\sigma),\label{gn:eq:S1-mean}\\
 \mathbb E_{n,\sigma}(S_3+12S_1^2)
      &=n\psi'''(2\sigma)+12n^2\psi'(2\sigma)^2.
          \label{gn:eq:fourth-insertion}
\end{align}
The complete paired deformation separates the two terms on the left
of the second identity and yields a stronger moment law.

\section{Exact finite Gaussian moments}
\label{gn:sec:gaussian}

\begin{theorem}[Gaussian moment matching through degree $n+1$]
\label{gn:thm:gaussian}
Let $n\ge1$, $\sigma>0$, and set
\begin{equation}\label{gn:eq:mu-v}
 \mu_\sigma=\psi'(2\sigma)>0,\qquad
 v_\sigma=\psi'''(2\sigma)>0,
 \qquad S=S_1.
\end{equation}
For every integer $0\le k\le n+1$,
\begin{equation}\label{gn:eq:S-moments}
 \boxed{\displaystyle
 \mathbb E_{n,\sigma}S^k
 =\sum_{j=0}^{\lfloor k/2\rfloor}
       \frac{k!}{2^jj!(k-2j)!}
       (n\mu_\sigma)^{k-2j}(v_\sigma/4)^j.}
\end{equation}
Equivalently, the statistic
\begin{equation}\label{gn:eq:X}
             X_{n,\sigma}
       =\frac{2S-2n\mu_\sigma}{\sqrt{v_\sigma}}
\end{equation}
has exactly the standard Gaussian moments through degree $n+1$:
\begin{equation}\label{gn:eq:Gaussian-moments}
 \mathbb E_{n,\sigma}X_{n,\sigma}^{k}
 =\begin{cases}
    0,&k\text{ odd},\\[2pt]
    (k-1)!!,&k\text{ even},
   \end{cases}
 \qquad 0\le k\le n+1,
\end{equation}
where the moment at degree zero is one.
\end{theorem}

\begin{proof}
Choose $k$ distinct pair labels in \eqref{gn:eq:polygamma-germ} and
extract the coefficient of $u_1^2\cdots u_k^2$.  Put $w_r=u_r^2$ and
compute in the quotient in which $w_r^2=0$.  On the left, all terms
with $u_r^4$ or higher disappear, so the required coefficient is
$\mathbb E_{n,\sigma}(2S)^k$.

For one pair interaction, Taylor expansion of its logarithm gives
\begin{align*}
 &\sum_{\varepsilon,\eta\in\{-1,1\}}
      \log\Gamma(2\sigma+\varepsilon u_r+\eta u_s)
                       -4\log\Gamma(2\sigma)\\
 &\hspace{12mm}\equiv
          2\mu_\sigma(w_r+w_s)+v_\sigma w_rw_s
                              \pmod{w_r^2,w_s^2}.
\end{align*}
The sum is even in each pair variable.  Each of the $n+1$ labels has
exactly $n$ partners, and no logarithmic summand depends on three
distinct labels.  Thus the logarithm of the complete right side reduces
to
\[
       2n\mu_\sigma\sum_{r=1}^kw_r
                  +v_\sigma\sum_{r<s\le k}w_rw_s.
\]
To obtain $w_1\cdots w_k$ from its exponential, partition the $k$
labels into singleton and pair blocks.  The number of partitions with
$j$ pairs is $k!/(2^jj!(k-2j)!)$.  Such a partition contributes
$(2n\mu_\sigma)^{k-2j}v_\sigma^j$.  Therefore
\[
 \mathbb E_{n,\sigma}(2S)^k
 =\sum_{j=0}^{\lfloor k/2\rfloor}
   \frac{k!}{2^jj!(k-2j)!}
       (2n\mu_\sigma)^{k-2j}v_\sigma^j.
\]
Division by $2^k$ proves \eqref{gn:eq:S-moments}.  These are the
moments of a Gaussian variable with mean $n\mu_\sigma$ and variance
$v_\sigma/4$.  Centering and scaling proves
\eqref{gn:eq:Gaussian-moments}.  The restriction $k\le n+1$ is exactly
the number of distinct pair parameters available in the proof.
\end{proof}

\begin{corollary}[Separated low-order polygamma integrals]
\label{gn:cor:variance}
For all $n\ge1$ and $\sigma>0$,
\begin{equation}\label{gn:eq:variance}
 \operatorname{Var}_{n,\sigma}(S_1)=\frac14\psi'''(2\sigma),
 \qquad
 \mathbb E_{n,\sigma}S_3=(n-3)\psi'''(2\sigma).
\end{equation}
At the quarter parameter,
\begin{equation}\label{gn:eq:quarter-statistics}
 \mathbb E_{n,1/4}S_1=\frac{n\pi^2}{2},\qquad
 \operatorname{Var}_{n,1/4}(S_1)=\frac{\pi^4}{4},\qquad
 \mathbb E_{n,1/4}S_3=(n-3)\pi^4.
\end{equation}
\end{corollary}

\begin{proof}
The degree-two instance of \eqref{gn:eq:S-moments} gives the variance.
Insert the resulting second moment into
\eqref{gn:eq:fourth-insertion} to obtain the expectation of $S_3$.
Finally, $\psi'(1/2)=\pi^2/2$ and $\psi'''(1/2)=\pi^4$ give the
quarter values \cite{dlmfpolygamma}.
\end{proof}

\subsection{Explicit one-dimensional identities}

At $n=1$ and $\sigma=1/4$, the normalizer is $4\pi^5$.  To display
several resulting evaluations compactly, write
\[
 \omega(x)=x\sinh(\pi x)|\Gamma(\tfrac14+ix/2)|^8,
 \qquad z_x=\tfrac14+ix/2.
\]
Every integral below is an ordinary absolutely convergent integral:
\begin{align}
 \int_0^\infty\omega(x)\dd x&=4\pi^5,
          \label{gn:eq:scalar-normalizer}\\
 \int_0^\infty\omega(x)\Re\psi'(z_x)\dd x&=2\pi^7,
          \label{gn:eq:scalar-trigamma}\\
 \int_0^\infty\omega(x)[\Re\psi'(z_x)]^2\dd x&=2\pi^9,
          \label{gn:eq:scalar-trigamma-square}\\
 \int_0^\infty\omega(x)\Re\psi'''(z_x)\dd x&=-8\pi^9.
          \label{gn:eq:scalar-third}
\end{align}
The negative sign in \eqref{gn:eq:scalar-third} is compatible with an
ordinary real integral: the real part of the complex-argument
polygamma function changes sign.

Unpaired parameter derivatives supply digamma identities as well.
Differentiating the $n=1$ version of
\eqref{gn:eq:determinant} once in one parameter and twice in two
distinct parameters, then setting all four parameters equal to $1/4$,
gives
\begin{align}
 \int_0^\infty\omega(x)\Re\psi(z_x)\dd x
    &=-4\pi^5(\gamma+3\log2),\label{gn:eq:scalar-psi}\\
 \int_0^\infty\omega(x)[\Re\psi(z_x)]^2\dd x
    &=4\pi^5\left((\gamma+3\log2)^2+\frac{\pi^2}{12}\right).
       \label{gn:eq:scalar-psi-square}
\end{align}
For completeness, the first logarithmic parameter derivative of the
Gamma quotient is $3\psi(1/2)-\psi(1)$, while the mixed logarithmic
derivative in distinct labels is $\psi'(1/2)-\psi'(1)$.  The integral
insertions are respectively $2\Re\psi(z_x)$ and
$4[\Re\psi(z_x)]^2$.  These factors give exactly
\eqref{gn:eq:scalar-psi}--\eqref{gn:eq:scalar-psi-square}.

\subsection{A distributional consequence of the exact identities}

The Gaussian matching is asserted only through the stated finite degree.
For fixed $n$ and $\sigma$ the statistic is bounded, since
\[
 |\Re\psi'(\sigma+iy)|\le\sum_{r=0}^\infty\frac1{(r+\sigma)^2}
                       =\psi'(\sigma).
\]
Thus it has a compactly supported distribution.  Its moment identities
nevertheless imply a Gaussian limit as the dimension increases.

\begin{corollary}[Gaussian limit for arbitrary positive parameter sequences]
\label{gn:cor:CLT}
For any sequence $\sigma_n>0$, the random variables
$X_{n,\sigma_n}$ of \eqref{gn:eq:X} converge in distribution to a
standard normal random variable as $n\to\infty$.
\end{corollary}

\begin{proof}
For every fixed integer $k$, the $k$th moment is exactly the standard
Gaussian moment once $n\ge k-1$, independently of $\sigma_n$.
The second moment equals one for every $n\ge1$, giving tightness.
For each fixed $k$, a higher even moment is uniformly bounded along
the tail of the sequence, so the $k$th powers are uniformly integrable.
Every subsequential weak limit therefore has all standard Gaussian
moments.  The normal distribution is moment determinate; for example
its even moments satisfy Carleman's condition.  Consequently every
subsequential limit is standard normal, proving the assertion.
\end{proof}

This consequence uses exact moment identities at each sufficiently
large finite dimension.  No uniform large-dimension estimate on Gamma
functions or on the parameter sequence is required.
